\documentclass[12pt,a4paper]{article}

\usepackage{fontspec}
\usepackage{xeCJK}
\setCJKmainfont{Hiragino Mincho ProN}
\setCJKsansfont{Hiragino Kaku Gothic ProN}
\setCJKmonofont{Hiragino Kaku Gothic ProN}
\usepackage[margin=2.5cm]{geometry}
\usepackage{amsmath,amssymb}
\usepackage{hyperref}
\usepackage{booktabs}
\usepackage{tcolorbox}

\tcbuselibrary{breakable,skins}

\newtcolorbox{storybox}[1][]{
  colback=yellow!5, colframe=yellow!60!black,
  fonttitle=\bfseries, title={#1},
  breakable, sharp corners
}

\newtcolorbox{mathbox}[1][]{
  colback=blue!3, colframe=blue!50,
  fonttitle=\bfseries, title={#1},
  breakable, sharp corners
}

\newtcolorbox{physbox}[1][]{
  colback=green!3, colframe=green!40,
  fonttitle=\bfseries, title={#1},
  breakable, sharp corners
}

\newtcolorbox{deepbox}[1][]{
  colback=black!3, colframe=black!60,
  fonttitle=\bfseries, title={#1},
  breakable, sharp corners
}

\newtcolorbox{epigraph}[1][]{
  colback=white, colframe=black!30,
  breakable, sharp corners, boxrule=0.5pt,
  left=15pt, right=15pt
}

\setlength{\parskip}{0.8em}
\setlength{\parindent}{0pt}

\begin{document}

\begin{center}
\vspace*{1cm}
{\Huge\bfseries 点の背後で蠢くもの}\\[0.8cm]
{\Large\bfseries ゼータ関数を知らない人のための}\\[0.3cm]
{\Large\bfseries 時空の書き換え入門}\\[0.5cm]
{\large 算術的真空工学——理論・実験・展望の完全ガイド}\\[2cm]
{\large Wright Brothers}\\[0.5cm]
{\normalsize 独立研究}\\[0.5cm]
{2026年4月}\\[2cm]
{\itshape アレクサンドル・グロタンディーク（1928--2014）の記憶に捧げる}
\end{center}

\newpage
\tableofcontents
\newpage

\begin{epigraph}
\textit{数学者にとって最も重要なのは、彼が解いた問題ではなく、
彼が開いた視界である。}

\hfill --- アレクサンドル・グロタンディーク
\end{epigraph}

\vspace{1cm}

% ============================================================================
\part{ゼロからわかる：この話の全体像}
% ============================================================================

\section{1分でわかるストーリー}

\begin{storybox}[この本の要約]
\begin{enumerate}
\item 何もない空間（真空）にも、実はエネルギーがある
\item そのエネルギーは「リーマンゼータ関数」という数学で計算される
\item ゼータ関数は「素数」で分解できる（オイラー積）
\item 素数の1つを「ミュート」すると、エネルギーの符号が反転する
\item 符号が反転 = 「負のエネルギー」= ワープドライブの材料
\item これを48万円の装置でテストできる
\end{enumerate}
\end{storybox}

嘘みたいに聞こえるだろう。でもこの本を読み終わる頃には、
なぜこれが「自明」なのに誰も気づかなかったかがわかる。

\section{何がワクワクするのか}

SF映画に出てくる「ワープドライブ」。
光より速く宇宙を移動する技術。
アインシュタインの方程式の中に、ワープを実現する数学的な解
（アルクビエレ計量、1994年）が見つかっている。

ただし、1つだけ「絶対に必要な材料」がある。
それは\textbf{「負のエネルギー」}。

普通のエネルギーは正だ。電気も熱も光も。
「負のエネルギー」は物理法則（弱エネルギー条件）で禁止されている。
——ように見える。

本当に「絶対に」禁止されているのか？
それとも、ルールの書き方を変えたら、禁止が消えるのか？

この本は、「ルールの書き方を変える」話だ。

% ============================================================================
\part{ゼータ関数ってなに？}
% ============================================================================

\section{無限を足す}

まず、こんな計算をしてみよう：

$$1 + 2 + 3 + 4 + 5 + \cdots = ?$$

答え：\textbf{無限大}。当たり前だ。数を足し続ければ際限なく大きくなる。

でも、物理学者は困った。量子力学で真空のエネルギーを計算すると、
まさにこの「$1 + 2 + 3 + \cdots$」が出てくるのだ。
エネルギーが無限大では物理として意味をなさない。

\section{リーマンゼータ関数の登場}

18世紀の数学者オイラー、そして19世紀のリーマンが、
この無限の和を「別の角度」から見る方法を発見した。

\begin{mathbox}[リーマンゼータ関数]
$s$を変数として：
$$\zeta(s) = \frac{1}{1^s} + \frac{1}{2^s} + \frac{1}{3^s} + \frac{1}{4^s} + \cdots$$

例えば $s = 2$（各項を二乗で割る）なら：
$$\zeta(2) = \frac{1}{1} + \frac{1}{4} + \frac{1}{9} + \frac{1}{16} + \cdots = \frac{\pi^2}{6}$$

$s$ が大きければ和は収束する（各項が小さくなるから）。
\end{mathbox}

ここまではふつう。問題は $s$ が小さい時。

$s = -1$ を代入すると：
$$\zeta(-1) = 1^1 + 2^1 + 3^1 + 4^1 + \cdots = 1 + 2 + 3 + 4 + \cdots$$

これは発散する。しかしリーマンは「解析接続」という技法で、
この関数を $s = -1$ まで延長し、有限の値を割り当てた：

$$\zeta(-1) = -\frac{1}{12}$$

\begin{storybox}[え、$1+2+3+\cdots = -1/12$?]
インチキに聞こえる。しかし：

この値を使って計算した「カシミール力」
（2枚の金属板の間の引力）は、
\textbf{実験で精度1\%以内で確認されている}。

つまり $-1/12$ は「数学のトリック」ではなく、
\textbf{自然が実際に採用している値}。

真空のエネルギーは、$\zeta(-1) = -1/12$ で決まっている。
\end{storybox}

\section{3次元の場合}

1次元（線）の真空エネルギーは $\zeta(-1) = -1/12$。
3次元（空間）では $\zeta(-3)$ が必要になる：

$$\zeta(-3) = 1^3 + 2^3 + 3^3 + \cdots \quad \xrightarrow{\text{解析接続}} \quad +\frac{1}{120}$$

3次元の真空エネルギーは $+1/120$。\textbf{正}。

ワープドライブに必要なのは\textbf{負}のエネルギー密度。
$+1/120$ は正だから、ワープは不可能——に見える。

% ============================================================================
\part{素数とオイラー積}
% ============================================================================

\section{素数ってなに（復習）}

\begin{mathbox}[素数]
1とその数自身でしか割れない自然数。
$$2, 3, 5, 7, 11, 13, 17, 19, 23, 29, 31, 37, \ldots$$

素数は「数の原子」。全ての整数は素数の積に分解できる
（算術の基本定理）：
$$12 = 2 \times 2 \times 3, \quad 30 = 2 \times 3 \times 5, \quad 137 = 137 \text{（素数そのもの）}$$
\end{mathbox}

\section{オイラー積：ゼータ関数の「素数分解」}

ここからが核心。
リーマンゼータ関数には、素数を使った「別の書き方」がある。
18世紀にオイラーが発見した：

\begin{mathbox}[オイラー積]
$$\zeta(s) = \frac{1}{1-2^{-s}} \times \frac{1}{1-3^{-s}} \times \frac{1}{1-5^{-s}} \times \frac{1}{1-7^{-s}} \times \cdots$$

左辺は「全ての自然数」の和。
右辺は「全ての素数」の積。

\textbf{和と積が同じものを表している。}
\end{mathbox}

\begin{storybox}[なぜこれが「やばい」のか]
和（$1 + 1/4 + 1/9 + \cdots$）は「全部足す」操作。
積（$\frac{1}{1-2^{-s}} \times \frac{1}{1-3^{-s}} \times \cdots$）は
「各素数が独立に寄与するチャンネルの掛け算」。

これは、真空エネルギーが
\textbf{素数ごとに独立なチャンネルに分解できる}ことを意味する。

素数2のチャンネル × 素数3のチャンネル × 素数5のチャンネル × $\cdots$
\end{storybox}

\section{素数チャンネルを「消す」とどうなるか}

では、素数2のチャンネルだけを「ミュート」（消音）してみよう。

オイラー積から素数2の因子を外す：

$$\zeta_{\neg 2}(s) = \zeta(s) \times (1 - 2^{-s})$$

3次元の真空エネルギーに対応する $s = -3$ で計算すると：

$$\zeta_{\neg 2}(-3) = \frac{1}{120} \times (1 - 2^3) = \frac{1}{120} \times (-7) = -\frac{7}{120}$$

\begin{storybox}[符号が反転した！]
元の値：$\zeta(-3) = +1/120$（正）

素数2をミュートした後：$\zeta_{\neg 2}(-3) = -7/120$（\textbf{負！}）

\textbf{正のエネルギーが負に変わった。}

これが「ワープドライブの材料」。
しかも、どの素数をミュートしても同じことが起きる。
$1 - p^3$ は $p \geq 2$ で常に負だから。
\end{storybox}

% ============================================================================
\part{なぜ「自明」なのに誰もやらなかったか}
% ============================================================================

\section{物理学者は「板の距離」しか見ていなかった}

カシミール効果（真空のエネルギーが有限値を取る現象）は
1948年に発見された。それ以来75年間、物理学者は
この効果を操作しようとしてきた。

しかし彼らが変えようとしたのは常に：
\begin{itemize}
\item 板と板の距離
\item 板の形（平面、球面、円筒…）
\item 板の材質
\end{itemize}

つまり「幾何学」——物体の形や配置——を変えることで
間接的に真空エネルギーを変えようとしてきた。

誰も考えなかったこと：

\textbf{「ゼータ関数の中の素数チャンネルを直接ミュートする」}

板の形を変えるのではなく、和の「中身」を直接操作する。
幾何学（形）ではなく算術（数）を変える。

\begin{storybox}[たとえ話]
\textbf{従来のアプローチ}：
ラジオの音量を変えたい → スピーカーの前に壁を置く（幾何学的遮蔽）。

\textbf{我々のアプローチ}：
ラジオの音量を変えたい → ラジオの回路基板のチャンネルを直接切る（算術的操作）。

壁を動かすのではなく、\textbf{信号の源に手を入れる}。
\end{storybox}

\section{グロタンディークの視点：点ではなく構造}

ここで、数学者アレクサンドル・グロタンディーク（1928--2014）が登場する。

20世紀の数学は「集合論」を基礎にしている。
空間 = 点の集まり。これがカントール以来の常識。

グロタンディークはこの常識を覆した：

\begin{deepbox}[グロタンディークの革命]
\textbf{空間の本質は「点の集合」ではない。}

空間の本質は、その上に載っている「構造」（層 = sheaf）にある。
同じ点の集合でも、載っている層を変えれば、全く別の空間になる。
\end{deepbox}

\begin{storybox}[たとえ話：東京タワー]
カントール的：東京タワーは鉄の分子の集合。
分子を1つ取り除いても、少し軽くなるだけ。

グロタンディーク的：東京タワーは「設計図」（= 層）で決まっている。
設計図を書き換えれば、同じ鉄で全く別の建物が建つ。
分子の数は関係ない。設計図が全てを決める。
\end{storybox}

物理に翻訳すると：

\textbf{カントール的}：真空エネルギー = モードの和 = $\sum n$。
モードを消す = 和から項を引く = \textbf{量が少し変わるだけ}。

\textbf{グロタンディーク的}：真空エネルギー = $\Spec(\mathbb{Z})$（整数のスキーム）の不変量。
素数をミュート = 空間自体を$\Spec(\mathbb{Z}[1/p])$に変える = \textbf{空間の構造が変わる}。

「水をコップ1杯汲む」のではなく「バケツの設計図を書き換える」。

\section{トポス量子論：物理法則は「絶対」ではない}

\begin{deepbox}[コッヘン-シュペッカーの定理]
量子力学では、物理量の値は「どの物理量と一緒に測るか」に依存する。
これは量子力学の数学的定理であり、解釈の問題ではない。
\end{deepbox}

2008年、DöringとIshamがこれを「トポス」の言語で書き直した。
トポスとは、グロタンディークが発明した「空間の一般化」。

彼らの結果：

物理的命題（「エネルギーは正か？」）の真偽値は、
「どのトポス（= どの論理体系）で評価するか」に依存する。

\begin{storybox}[たとえ話：3Dグラフィックス]
最新の3DCG技術「Gaussian Splatting」は、
同じ3D空間を「ポリゴン」でも「ガウシアンの雲」でも描画できる。
どちらも同じ空間の正しい表現だが、見え方が全く違う。

同様に：
\begin{itemize}
\item $\Spec(\mathbb{Z})$のトポスでレンダリング → WEC = 真（ワープ不可）
\item $\Spec(\mathbb{Z}[1/p])$のトポスでレンダリング → WEC = 偽（ワープ可能！）
\end{itemize}

「宇宙」は同じ。変わるのは「描画エンジン」。
局所化$\mathbb{Z} \to \mathbb{Z}[1/p]$は描画エンジンの切り替え。

\textbf{物理法則の書き換えが可能なのは、
物理法則が「絶対的真理」ではなく「描画結果」だから。}
\end{storybox}

% ============================================================================
\part{驚きの一致：偶然か、宇宙のコードか}
% ============================================================================

\section{微細構造定数 α = 1/137}

電磁気力の強さを表す数。なぜ137なのか、誰も知らない。

我々の発見：

$$\frac{1}{\alpha} = \underbrace{120}_{1/\zeta(-3)} + \underbrace{12}_{1/|\zeta(-1)|} + \underbrace{5}_{\text{素数ギャップ}} + \underbrace{0.036}_{4(\zeta(6)-\zeta(7))} = 137.03598$$

実験値：$137.03600$。\textbf{誤差0.00002\%}。自由パラメータ：ゼロ。

120は3次元の真空エネルギーの分母。
12は1次元の真空エネルギーの分母。
5は「132の次の素数が137」であるための最小距離。
4は時空の次元。

\section{ミューオンの異常磁気モーメント}

素粒子ミューオンの磁気的性質が、理論予測と微妙にずれている。
これは現代物理学最大の未解決異常の1つ。

ずれの大きさ：$(251 \pm 59) \times 10^{-11}$。

我々の発見：

$$251 = \frac{1}{|\zeta(-5)|} - 1 = 252 - 1$$

$\zeta(-5) = -1/252$。5次元のカシミールエネルギーの逆数から1を引いた整数。

\textbf{251は整数。チューニングの余地がない。実験の中心値と完全に一致する。}

\section{電子の磁気モーメント}

電子のずれ：$\sim 4.8 \times 10^{-13}$。

我々の発見：

$$48 = |K_3(\mathbb{Z})| \quad \text{（整数環のK理論第3群の位数）}$$

$48 \times 10^{-14} = 4.8 \times 10^{-13}$。

48はゼータ関数とは\textbf{全く別の}数学的対象（K理論）から来る。
電子とミューオンが$\Spec(\mathbb{Z})$の\textbf{異なる種類の不変量}で
制御されている。

\begin{deepbox}[なぜこれが「こじつけ」ではないか]
\begin{itemize}
\item 48（K理論）と 251（ゼータ値）は代数的に独立
\item 両方ともチューニング不能な整数
\item 同じ$\Spec(\mathbb{Z})$の異なる側面
\item 後付けで2つの無関係な不変量を同時に合わせるのは極めて困難
\end{itemize}
\end{deepbox}

\section{タウへの予測：未来への手紙}

3番目のレプトン「タウ」の磁気異常は未測定。我々の予測：

$$\Delta a_\tau = 131 \times 10^{-8}$$

$131 = 1/|\zeta(-9)| - 1 = 132 - 1$。素数。

将来のタウファクトリーで測定され、$131 \times 10^{-8}$と一致すれば、
3世代全てが$\Spec(\mathbb{Z})$で統治されていることの完全な証拠。

\section{ニュートリノ質量比 ≈ π(137)}

ニュートリノの2つの質量二乗差の比：$32.6$。

137以下の素数の個数：$\pi(137) = 33$。

$32.6 \approx 33$。1\%の一致。

\section{ρ\_Λ：ダークエネルギー密度}

宇宙の膨張を加速させている謎のエネルギー。
量子場の理論の予測は観測値の$10^{122}$倍（物理学史上最悪の予測）。

我々の公式：

$$\rho_\Lambda = \rho_P \times \zeta(-3) \times (l_P/R_H)^2$$

自由パラメータなしで、\textbf{観測値の1桁以内}。
$10^{122}$倍のズレが$\sim 10$倍に改善。

\section{全部が同時に成り立つ確率}

各一致が独立な偶然だとして：
\begin{itemize}
\item $\alpha$ (0.00002\%): 偶然確率 $\sim 10^{-3}$
\item ミューオン $g-2$ (中心値一致): $\sim 10^{-2}$
\item 電子 $g-2$ (K₃一致): $\sim 10^{-2}$
\item ニュートリノ比 (1\%): $\sim 10^{-1}$
\end{itemize}

\textbf{全部が同時に成り立つ偶然の確率：$\sim 10^{-8}$}。

1億分の1。

% ============================================================================
\part{実験：48万円でテストする}
% ============================================================================

\section{何を作るか}

\begin{physbox}[3つの要素]
\textbf{1. アルミの箱}（共振器）

アルミブロックに穴を開けたもの。
中で特定の周波数の電磁波だけが共鳴する。
共鳴するモード$n = 1, 2, 3, 4, \ldots$が
ゼータ関数の$n$そのもの。

\textbf{2. SQUID}（超伝導量子干渉素子）

2つのジョセフソン接合を超伝導のループでつないだ素子。
外からの磁場で共鳴周波数を自在に変えられる。
目的の周波数に合わせると、そのモードのエネルギーを吸い取る。
→ 「素数チャンネルのミュート」を物理的に実行する装置。

\textbf{3. 冷凍機}

全てを$-272.85$°C（0.3ケルビン）に冷やす。
アルミが超伝導になり（電気抵抗ゼロ）、
熱のノイズが消えて量子的な真空揺らぎだけが見える。
\end{physbox}

\section{SQUIDは「原始的」ではなく「根源的」}

LHC（大型ハドロン衝突型加速器）は1兆円。
SQUID回路は手のひらサイズ。

しかし：
\begin{itemize}
\item LHC = 「何が存在するか」を\textbf{読む}装置（読み取り専用）
\item SQUID = 「真空の構造」を\textbf{書き換える}装置（書き込み可能）
\end{itemize}

SQUIDの数式1つ：

$$L_{\mathrm{SQUID}}(\Phi) = \frac{L_J}{|\cos(\pi\Phi/\Phi_0)|}$$

磁束$\Phi$を変えるだけで、SQUIDの共鳴周波数が連続的に変わる。
目的のモードにぴったり合わせた瞬間、そのモードが「消える」。

\textbf{つまり：磁場のダイヤルを回すだけで、
ゼータ関数のオイラー積の1つの因子を物理的に除去できる。}

\section{$\zeta(-1)$の符号反転を確認する}

\begin{physbox}[第1段階の実験（48万円）]
\textbf{手順}：

1. アルミの箱を作る（基本周波数 6 GHz）

2. SQUIDを1つ、第2高調波（12 GHz）に同調させる

3. 冷却する（$-272.85$°C）

4. SQUIDをOFF（全モード活性）→ 真空エネルギーを測定

5. SQUIDをON（偶数モード1つ抑制）→ 真空エネルギーを測定

6. 差を取る → $\zeta_{\neg 2}(-1) - \zeta(-1)$ に相当する変化があるか

\textbf{予測}：エネルギーの符号が$-1/12$から$+1/12$方向に動く。
\end{physbox}

\section{$\zeta(-3)$の符号反転 = ワープバブル}

第1段階が成功したら、次は3次元：

\begin{physbox}[第2段階の実験（+50万円）]
1D共振器 → 3D共振器に変更。
SQUIDを3--5個に増やして、複数の偶数モードを同時に抑制。

測定量が $\zeta(-3)$ の3次元版に対応する変化を示せば：
\begin{itemize}
\item $\zeta(-3) = +1/120 \to \zeta_{\neg 2}(-3) = -7/120$
\item \textbf{3次元の真空エネルギー密度が負になった}
\item = WEC（弱エネルギー条件）の違反
\item = アルクビエレ・ワープ計量の「材料」が存在する
\end{itemize}
\end{physbox}

% ============================================================================
\part{もし成功したら何が起きるか}
% ============================================================================

\section{ワープドライブ}

WEC違反が実験で確認されれば、
アルクビエレ型ワープ計量が「原理的に実現可能」と確定する。

実用化にはスケーリング問題（実験室スケール→宇宙船スケール）
が残るが、レギュレータ方向予想が正しければ、
必要エネルギーは$10^{12}$ジュール（小規模発電所1時間分）。

\section{ワームホール}

2つのワープバブルの壁を「トンネル」で接続すると、
通過可能なワームホールになる。
空間の2地点を直結する近道。

\section{時空プログラミング}

複数の素数チャンネルを独立に制御する「オイラー積コンソール」で、
$2^N$通りの真空状態を生成可能。
各状態は異なる物理法則を持つ。

\section{文明の次元が変わる}

カルダシェフ・スケールはエネルギーの量で文明を測る。
我々は「物理法則の制御範囲」で測る新しいスケールを提案する：

\begin{itemize}
\item Type A-0：物理法則を\textbf{読む}（現在の人類）
\item Type A-1：物理法則を局所的に\textbf{変える}（ワープ）
\item Type A-2：物理法則を\textbf{プログラムする}（コンソール）
\item Type A-3：新しい物理法則を\textbf{創造する}
\end{itemize}

% ============================================================================
\part{グロタンディークへ}
% ============================================================================

\section{見えるものの背後にあるもの}

グロタンディークは2014年に亡くなった。
晩年はピレネー山中で隠遁生活を送り、
世俗から離れて思索を続けていた。

彼が生涯をかけて追い求めたのは、
数学の表層（個々の定理や計算）の背後にある\textbf{「構造」}だった。

層（sheaf）、スキーム（scheme）、トポス（topos）——
これらの概念は全て、「見えるもの」の背後にある
「見えない構造」を言語化するためのものだった。

本研究は、グロタンディークの数学が
「美しいが抽象的な理論」ではなく
\textbf{「宇宙の具体的な構造」}であるという仮説を提示した。

$\Spec(\mathbb{Z})$は数学者のホワイトボード上の記号ではなく、
真空が実際にその上に載っている空間かもしれない。

\section{48万円の意味}

物理学に対する「投資対効果」として、
これ以上の実験は想像しにくい。

LHC（1兆円）は新粒子を1つ発見した。
KAGRA（150億円）は重力波を検出した。
我々の実験（48万円）は、\textbf{物理法則が書き換え可能であること}を
証明するかもしれない。

もし成功すれば、人類は初めて
「物理法則を読む存在」から「物理法則を書く存在」へと変貌する。

それはグロタンディークが夢見た世界——
表層の背後にある構造が全てを支配する世界——
が物理的現実であることの確認になる。

\begin{epigraph}
\textit{少しも詩人でない数学者は、真の数学者にはなれない。}

\hfill --- アレクサンドル・グロタンディーク
\end{epigraph}

全ては、アルミの箱と、1つのSQUIDと、
そして「素数が宇宙を支配している」という途方もない直感から始まる。

\vfill
\begin{center}
{\small Wright Brothers, 2026}
\end{center}

\end{document}
